\documentclass[12pt]{article}
\usepackage[utf8]{inputenc}
\usepackage[T1]{fontenc}
\usepackage[margin=1in]{geometry}
\usepackage{setspace}
\usepackage{xcolor}
\usepackage[colorlinks=true, urlcolor=blue!60!black, linkcolor=blue!60!black]{hyperref}
\setlength{\parindent}{2em}
\setlength{\parskip}{1.25em}
\renewcommand{\baselinestretch}{1.0}
\begin{document}
\begin{center}
{\LARGE \textit{United States v. Belmont}, 301 U.S. 324 (1937)\par}
\vspace{0.8em}
{\large Case Brief --- Four-Page Edition\par}
\end{center}
\vspace{1em}

\section{Memory Jogger}

The Court unanimously reversed the dismissal of the United States' suit to recover a New York bank deposit assigned to it by the Soviet Union in the Litvinov Assignment; Justice Sutherland's majority held that a sole executive agreement made incident to recognition needs no Senate consent and overrides contrary state law and policy (\href{https://supreme.justia.com/cases/federal/us/301/324/}{Justia}).

\section{Detailed Case Facts}

Before 1918, the Russian corporation Petrograd Metal Works deposited money with August Belmont, a private banker in New York City. In 1918, after the Bolshevik Revolution, the Soviet Government nationalized the corporation and expropriated all its property wherever situated, including the New York deposit.

The United States withheld recognition of the Soviet regime for sixteen years. On November 16, 1933, President Franklin D. Roosevelt recognized the Soviet Government. In an exchange of diplomatic correspondence with Soviet Foreign Commissar Maxim Litvinov---the ``Litvinov Assignment''---never submitted to the Senate, the Soviet Government assigned to the United States all amounts due to it from American nationals, including the Petrograd Metal Works deposit, so that the United States could assemble a fund to settle Americans' claims against Russia.

\section{Procedural History}

After Belmont's death, the United States demanded the deposit from his executors and sued in the Southern District of New York. The district court dismissed for failure to state a cause of action, and the Second Circuit affirmed, reasoning that the deposit's situs was New York and that enforcing the Soviet decree would give effect to a confiscation contrary to New York public policy. The Supreme Court granted certiorari, heard argument on March 4, 1937, and decided the case on May 3, 1937.

\section{Judicial Votes}

Reversed, \textbf{9--0} in the judgment.
\begin{itemize}
\item \textbf{Majority:} Justice Sutherland, joined by Chief Justice Hughes and Justices Van Devanter, McReynolds, Butler, and Roberts.
\item \textbf{Concurring in the judgment:} Justice Stone, joined by Justices Brandeis and Cardozo.
\item \textbf{Dissents:} None.
\end{itemize}

\section{Holding}

Recognition retroactively validated the Soviet Government's acts from its inception, barring American courts from questioning its taking of its own corporation's property. The President, as the organ of external relations, could make the Litvinov Assignment as part of recognition without Senate participation. A valid sole executive agreement carries federal supremacy in foreign affairs and overrides conflicting state laws and policies, including New York's policy against confiscation.

\section{Analysis of Opinions}

\textbf{Majority (Sutherland).} \emph{Recognition:} under \textit{Underhill v. Hernandez} and \textit{Oetjen v. Central Leather}, recognition validates a foreign government's acts from the beginning, and courts may not judge its expropriation of its own nationals' property. \emph{Executive agreements:} drawing on his theory of inherent external sovereignty in \textit{Curtiss-Wright} (1936), Sutherland reasoned that the power to recognize a government includes the power to settle claims obstructing relations. A ``treaty'' requires two-thirds of the Senate, but protocols, \textit{modi vivendi}, postal conventions, and claims assignments fall within the President's competence. \emph{Supremacy:} although Article VI names only ``Treaties,'' valid executive agreements have the same preemptive force; in foreign affairs ``state lines disappear,'' and for those purposes ``the State of New York does not exist.'' \emph{Extraterritoriality:} the Constitution has no extraterritorial operation except as to citizens; the Russian corporation's redress lay with its own government, and the executors were mere stakeholders. Rights of any domestic creditors were left open.

\textbf{Stone (concurring in the judgment).} A State ordinarily may refuse to enforce foreign confiscations of property within its borders, but New York had no policy letting a mere stakeholder resist payment to a recognized sovereign's assignee when no local creditors had intervened. He refused to hold that a sole executive agreement overrides state law: neither document purported to alter New York law or give the United States greater rights than the Soviet Union had.

\textbf{Aftermath.} \textit{United States v. Pink} (1942) expanded \textit{Belmont} over Stone's dissent; \textit{Dames \& Moore} and \textit{Garamendi} relied on it; \textit{Medell\'in v. Texas} confined it to a ``narrow set of circumstances''; and \textit{Zivotofsky v. Kerry} cited it for the exclusive recognition power.

\section{Examples: Future Applications}

\begin{enumerate}
\item \textbf{Same-Side (Preempted): Settlement Incident to Recognition.} A sole executive agreement accompanying recognition of a new government extinguishes U.S. claims against its airline in exchange for a compensation fund; a contrary New York statute yields.
\item \textbf{Same-Side (Collection Allowed): Assigned Bank Claim.} A new government assigns its nationalized central bank's claims to the United States; a State's anti-confiscation policy cannot block collection.
\item \textbf{Opposite-Side (Not Preemptive): Directing State Criminal Courts.} An unratified agreement promising that state courts will vacate convictions, without implementing legislation, fails under \textit{Medell\'in}.
\item \textbf{Opposite-Side (Statute Prevails): Waiving Statutory Duties.} An executive agreement promising duty-free steel contrary to a federal statute cannot displace an Act of Congress; \textit{Belmont} addressed only state law.
\item \textbf{Fence-Sitter: Ending Tort Suits Without a Fund.} Terminating state-court suits against a long-recognized government's company for trade concessions, without a substitute remedy or recognition, sits at the edge of the ``narrow set of circumstances'' \textit{Medell\'in} allowed.
\end{enumerate}

\section{Critique}

\textbf{Text.} Article II requires two-thirds Senate consent for treaties, and Article VI makes supreme only the Constitution, federal statutes, and ``Treaties.'' Letting a sole executive agreement preempt state law lets the President make supreme domestic law alone, without bicameralism or a Senate supermajority.

\textbf{Overbreadth and asymmetry.} Stone's concurrence exposed the ``state lines disappear'' rhetoric as unnecessary maximalism, since States govern local property, estates, and banking; \textit{Medell\'in} later vindicated his caution. And because a sole executive agreement cannot override a federal statute, such agreements occupy an awkward intermediate tier of the supremacy hierarchy.

\section{Key Quotations}

\begin{quote}
``Plainly, the external powers of the United States are to be exercised without regard to state laws or policies.''
\end{quote}
\noindent --- Justice Sutherland, for the Court, 301 U.S. at 331.

\begin{quote}
``There are many such compacts, of which a protocol, a modus vivendi, a postal convention, and agreements like that now under consideration are illustrations.''
\end{quote}
\noindent --- Justice Sutherland, for the Court, 301 U.S. at 330--331.

\begin{quote}
``... our Constitution, laws and policies have no extraterritorial operation, unless in respect of our own citizens. ... What another country has done in the way of taking over property of its nationals, and especially of its corporations, is not a matter for judicial consideration here. Such nationals must look to their own government for any redress to which they may be entitled.''
\end{quote}
\noindent --- Justice Sutherland, for the Court, 301 U.S. at 332.

\end{document}
